\chapter*{译\ 者\ 序}
\phantomsection
\addcontentsline{toc}{chapter}{译者序}
\markboth{译者序}{译者序}

这本书到我们手里的时候，是一份日文打印稿和一只硬盘。打印稿有两百多页，页边写满了作者自己的批注，有些批注比正文还长。硬盘里收着同一部稿子的电子版，另加三百来张翻拍照片，编号混杂，文件名里夹着俄文、日文和数字。

原稿是分章交来的。孙笑川译完一章，李文亚再核这一章里的俄文档案引文，两边的进度始终错开一章。头三个月做的全是最琐碎的活：把照片按 \texttt{ФВ} 编号排好，再把稿子里提到而硬盘里缺着的照片挑出来，列成一张缺件表。表列出来一共缺四十一处，后来补齐了三十九处，还剩两处一直没找到：一处是背面某片低地的原片，作者说他自己也没见过；另一处是某张表的第七张特写，据说拍过，翻遍硬盘也没有。正文里凡是涉及这两处的地方，我们都照原稿译出，没有替作者含糊过去。

体例上只有三条要交代。第一，本书引及《战争的雷霆》与《月之都——〈广寒宫残简〉辑佚与考释》时不注页码，只注书名与相关章节，因为各书版次不同、页码会变动，章节篇目不会。第二，原著正文一律用「笔者」，序言与结语用「我」。日文初稿原本通篇写作「私」，我们建议他按体例区分开来，作者答复说这样确实更清楚。第三，译文不另加注，也不替作者补充材料，专名一律以能与另外两部书对齐的译法为准。

还有一件事不在体例之内，却也必须放在这里说清。这部稿子由两人合译，而此前我们曾为另一部稿子闹翻过，此后各自翻译自己的部分，连校样都不互相过目。这次勉强算作例外，只因材料横跨三部书，一个人实在译不下来。译文里若出现前后口径不太一样的地方，我们一律原样保留。读者读到互相抵牾之处，那多半是我们各译各的结果，未必是哪一边译错了。

\begin{flushright}
{\small 孙笑川}\quad{\small 李文亚}\\[2mm]
{\footnotesize\color{pencil} 二〇二四年十一月，于北京}
\end{flushright}

\clearpage
% ===================================================================
\chapter*{著者中文版序言}
\phantomsection
\addcontentsline{toc}{chapter}{著者中文版序言}
\markboth{著者中文版序言}{著者中文版序言}

我没去过月面。这一点先讲清楚。我见过的只有照片、一卷缩微胶片，还有巴兰尼科夫教授书里那几张重绘图。书里凡是讲「背面长什么样」的地方，读者都可以先打个折，因为那些地方我自己也没踩上去看过。

校样是这样到我手上的。二〇二四年二月，事务所收到一只包裹，寄件人写着一家财团的名字，品名栏填的是「同人志样本」。这类邮件每年都有几十件，我照例拆开，里面是一叠装订成册的俄文校样，封面印着丛书的式样。第一页夹着一张手写便条，上面只有一句「您可能想看看这个」，没有署名；翻到背面还有一行小字：「地址是九〇年代留下的，一直没注销。」

那天下午我把图翻了一遍。俄文一个字也看不懂，我就把正文一页一页拍下来丢进翻译软件，翻出来的日文半通不通，只能对着图猜。就这样看了三天，看到一半才明白，我读的不是别人编的东西。又过了三天，我订了去莫斯科的机票。

二〇二四年三月我到了莫斯科，专门去看那批实物和照片。去之前带了三本笔记本，回来只用掉半本。东西其实不少，只是我大多认不出来：照片上的地形我看得明白，照片底下那行俄文编号和注释我看不懂。等那一叠校样全部请人译出来，已经是两个月以后的事。

还有一件更让人头疼的事。那些物件当年似乎上过新闻，报纸登过消息，展览办过，大人还排过队。这些都是别人告诉我的，我自己一点印象也没有。按年纪算，那时候我大概也在电视机前坐过，现在却一个字也想不起来，脑子里只剩一句「好像确实上去过一次」。另一头就更麻烦了：这些零碎记忆后来全进了我的作品里。我做了几十年创作，把月亮画成过好几种样子，画里那座城有名字、有居民，也有一套自己的规矩，不少读者都认得。只是大家认得它，是把它当成一部娱乐作品里的设定。我当年做的时候，也只觉得这样写很有意思，现在回头翻档案才明白，那大概是我把记得的片段换了个更好看、更容易让人听下去的讲法重新编了一遍。我做了三十年游戏，画过很多次月亮，没有一张画对过。问题出在我从一开始就搞错了它在哪儿、有多大、里面装了什么。

他们留在月面上时确实不会变化，一落地倒立刻开始变了，而且变得很快。第一次看校样时我担心这本书会被读成一篇寓言，其实我没往那边写。材料摆在桌上长什么样，我只照着量出来、抄下来。书里每一个数字都有出处，对不上的都标出来了。前半部分我有图有表核得动，写得细；后半部分材料散，只能一步步往下推。读者若觉得前面慢、后面赶，原因就在这里。有一处我说得比手头材料能支持的还要肯定，这一点放在结语里认账。

\begin{flushright}
{\small 太田顺也}\\[2mm]
{\footnotesize\color{pencil} 二〇二四年十二月，于东京}
\end{flushright}

\clearpage
